\begin{table}[htbp]
\centering
\caption{Dose-Response: Effects by Item Quantity}
\label{tab:dose_response}
\begin{adjustbox}{max width=\textwidth}
\begin{threeparttable}
\small
\begin{tabular}{lcccc}
\toprule
 & \multicolumn{2}{c}{Log prices} & \multicolumn{2}{c}{Log firms} \\
\cmidrule(lr){2-3} \cmidrule(lr){4-5}
 & Base & PBU FE & Base & PBU FE \\
\midrule
\multicolumn{5}{l}{\textit{Panel A: High quantity (above median)}} \\
$g65 \times Pre$ & -0.0963*** & -0.1037*** & 0.1248*** & 0.1393*** \\
 & (0.0153) & (0.0149) & (0.0094) & (0.0093) \\
Observations & 348,363 & 348,363 & 402,194 & 402,194 \\
\midrule
\multicolumn{5}{l}{\textit{Panel B: Low quantity (below median)}} \\
$g65 \times Pre$ & -0.1220*** & -0.1189*** & 0.0251*** & 0.0268*** \\
 & (0.0129) & (0.0122) & (0.0086) & (0.0084) \\
Observations & 301,351 & 301,351 & 371,069 & 371,069 \\
\bottomrule
\end{tabular}
\begin{tablenotes}
\small
\item \textit{Notes:} 18-month window. The sample is split at the median of log quantity (3.18). High-quantity items are more likely to attract large non-SME suppliers, so the SME restriction is more binding for these items. This split aligns with the causal forest finding that item quantity is the dominant source of treatment effect heterogeneity (variable importance = 0.486). SE clustered at item level. * $p<0.10$, ** $p<0.05$, *** $p<0.01$.
\end{tablenotes}
\end{threeparttable}
\end{adjustbox}
\end{table}
